% Part: first-order-logic
% Chapter: axiomatic-deduction
% Section: identity

\documentclass[../../../include/open-logic-section]{subfiles}

\begin{document}

\olfileid[fr]{fol}{axd}{ide}

\olsection{\usetoken{P}{derivation} avec l'\usetoken{s}{identity}}

Pour prendre en charge~$\eq$ dans les !!{derivation}s, il suffit
d'ajouter de nouveaux schémas d'axiomes. La définition d'une
!!{derivation} et celle de~$\Proves$ restent inchangées; nous
admettons simplement ces axiomes supplémentaires.

\begin{defn}[Axiomes de l'!!{identity}]
\begin{align}
\ollabel{ax:id1} & \eq[t][t], \\
\ollabel{ax:id2} & \eq[t_1][t_2] \lif (!B(t_1) \lif !B(t_2)),
\end{align}
pour tous termes clos $t$, $t_1$ et~$t_2$.
\end{defn}

\begin{prop}\ollabel{prop:sound}
L'\olref{ax:id1} et l'\olref{ax:id2} sont des axiomes valides.
\end{prop}

\begin{proof}
Exercice.
\end{proof}

\begin{prob}
Démontrer la \olref[fol][axd][ide]{prop:sound}.
\end{prob}

\begin{prop}\ollabel{prop:iden1}
$\Gamma \Proves \eq[t][t]$ pour tout terme clos~$t$ et tout
ensemble~$\Gamma$.
\end{prop}

\begin{prop}\ollabel{prop:iden2}
Si $\Gamma \Proves !A(t_1)$ et $\Gamma \Proves \eq[t_1][t_2]$, où
$t_1$ et $t_2$ sont des termes clos, alors
$\Gamma \Proves !A(t_2)$.
\end{prop}

\begin{proof}
La !!{formula}
\[
(\eq[t_1][t_2] \lif (!A(t_1) \lif !A(t_2)))
\]
est une instance de l'\olref{ax:id2}. La conclusion découle de deux
applications de~\MP.
\end{proof}

\begin{editorial}
La définition source réserve les deux axiomes aux termes clos, mais
les deux propositions suivantes parlent ensuite de termes quelconques.
Le qualificatif « clos » y est rétabli afin que leurs conclusions
soient bien des instances des axiomes énoncés.
\end{editorial}

\end{document}
